% Paper B: h=1 CM field selection rule in Kanno-Watari W=0 vacua
% Status: CONJECTURAL PREPRINT — evidence-based, not a rigorous theorem
% Verified citations marked [V]; unverified marked [CITE-NEEDED]

\documentclass[12pt,a4paper]{article}

\usepackage{amsmath,amssymb,amsthm}
\usepackage{hyperref}
\hypersetup{hidelinks}
\usepackage{geometry}
\geometry{margin=2.5cm}
\usepackage[numbers,sort&compress]{natbib}
\usepackage{microtype}

\theoremstyle{plain}
\newtheorem{theorem}{Theorem}[section]
\newtheorem{conjecture}[theorem]{Conjecture}
\newtheorem{proposition}[theorem]{Proposition}
\newtheorem{lemma}[theorem]{Lemma}
\newtheorem{corollary}[theorem]{Corollary}

\theoremstyle{definition}
\newtheorem{definition}[theorem]{Definition}
\newtheorem{example}[theorem]{Example}
\newtheorem{remark}[theorem]{Remark}

\newcommand{\Q}{\mathbb{Q}}
\newcommand{\Z}{\mathbb{Z}}
\newcommand{\C}{\mathbb{C}}
\newcommand{\R}{\mathbb{R}}
\newcommand{\OK}{\mathcal{O}_K}
\newcommand{\HCF}{\mathrm{HCF}}
\newcommand{\Gal}{\mathrm{Gal}}
\newcommand{\End}{\mathrm{End}}
\newcommand{\Hom}{\mathrm{Hom}}
\newcommand{\rank}{\mathrm{rank}}
\newcommand{\disc}{\mathrm{disc}}

\title{Class Number One as a Vacuum Selection Rule:\\
CM Field Constraints in Kanno--Watari $W=0$ Flux Compactifications}

\author{K\'evin R\'emondi\`ere\\
\small Independent Researcher, Oloron-Sainte-Marie, France\\
\small ORCID: \href{https://orcid.org/0009-0008-2443-7166}{0009-0008-2443-7166}\\
\small \texttt{kevin.remondiere@gmail.com}}

\date{May 2026 \\ \small\textit{Preprint — conjectural bridge, not a theorem}}

\begin{document}

\maketitle

\begin{abstract}
We propose and motivate a vacuum selection conjecture for type IIB flux
compactifications on $K3 \times K3$ following the Kanno--Watari $W=0$
stabilization framework \cite{KannoWatari2024}. Specifically, we conjecture that
when both $K3$ factors carry complex multiplication (CM) by the same imaginary
quadratic field $K = \Q(\sqrt{-D})$, the resulting moduli vacuum is
\emph{isolated} (a single vacuum rather than a continuous family or Galois
orbit of vacua) \emph{if and only if} $K$ has class number $h(-D) = 1$.
The ``only if'' direction is supported by classical algebraic number theory:
when $h(-D) > 1$, the relevant special values of the L-function $L(f, 2)$
associated to the $K3$ newform escape $\Q$ and instead live in the Hilbert
class field of $K$, generating a Galois orbit of cardinality $h(-D)$,
which we interpret as $h(-D)$ physically inequivalent vacua. The ``if''
direction is corroborated by a proved hierarchy of L-value rationality
(M142, ECI framework) for each of the nine Heegner discriminants. This
conjecture constitutes a testable arithmetic bridge between string vacuum
geometry and algebraic number theory, accessible to verification via the
LMFDB and PARI/GP.
\end{abstract}

\tableofcontents

\newpage

%===========================================================================
\section{Introduction}
\label{sec:intro}
%===========================================================================

The problem of vacuum selection in string theory --- understanding why the
observed universe occupies one point in an astronomically large landscape ---
is one of the central open questions in theoretical physics. Flux
compactifications offer a concrete framework for discretizing the landscape,
but explicit arithmetic criteria that single out individual vacua remain rare.

Kanno and Watari \cite{KannoWatari2024} studied type IIB flux
compactifications on $K3 \times K3$ with the $W=0$ condition imposed as a
Hodge-theoretic constraint on the choice of flux. Their construction realizes
exact moduli stabilization for K3 surfaces whose period lattices satisfy
specific arithmetic conditions, notably the presence of complex multiplication.
This produces supersymmetric $\mathrm{AdS}_4$ (or Minkowski) vacua whose
moduli are fixed by the CM structure rather than by a generic flux search.

A classical result in arithmetic geometry, going back to Gauss and made
precise by Kronecker--Weber and class field theory, is that imaginary
quadratic fields $K = \Q(\sqrt{-D})$ split into two regimes according to
their class number $h(-D)$:
\begin{itemize}
  \item $h(-D) = 1$ (the nine Heegner fields, $D \in \{1, 2, 3, 7, 11,
    19, 43, 67, 163\}$): $\OK$ is a principal ideal domain; the Hilbert
    class field of $K$ coincides with $K$ itself; every CM value is
    already rational over $K$.
  \item $h(-D) > 1$: the Hilbert class field $\HCF(K)$ is a non-trivial
    extension of degree $h(-D)$ over $K$; CM values of modular functions
    live in $\HCF(K) \supsetneq K$, producing algebraic conjugates.
\end{itemize}

Moore's ``Arithmetic and Attractors'' \cite{Moore1998} established a deep
connection between BPS black hole attractor equations and CM points of
moduli spaces, showing that enhanced symmetry and arithmetic at CM loci
have physical content. The BIKMSV analysis \cite{BIKMSV1996} connected
singularity enhancement to rational points in moduli space. Together these
suggest that the rationality or irrationality of certain period-derived
quantities is physically meaningful.

The present paper proposes a precise version of this intuition adapted to
the KW setting: \textbf{Conjecture P1} asserts that the class number $h(-D)=1$
condition is exactly the arithmetic criterion distinguishing isolated vacua
from Galois families of vacua in the KW $W=0$ stabilization.

\paragraph{Organization.}
Section~\ref{sec:kw} reviews the KW framework and the role of CM K3 surfaces.
Section~\ref{sec:conjecture} states Conjecture P1 precisely.
Section~\ref{sec:evidence} presents the evidence from L-value rationality
(M142 hierarchy).
Section~\ref{sec:hcf} gives the Hilbert class field interpretation.
Section~\ref{sec:gap} discusses the gap between heuristic and rigorous proof.
Section~\ref{sec:open} lists open problems and computational tests.

%===========================================================================
\section{Kanno--Watari $W=0$ Stabilization and CM K3 Surfaces}
\label{sec:kw}
%===========================================================================

\subsection{Setup and notation}

We work with type IIB string theory compactified on $X = K3_1 \times K3_2$,
with $\mathcal{N}=2$ supersymmetry in $d=4$. The Gukov--Vafa--Witten
superpotential is
\[
  W = \int_X G_3 \wedge \Omega,
\]
where $G_3 = F_3 - \tau H_3$ is the complexified three-form flux and
$\Omega$ is the holomorphic volume form. The $W=0$ condition requires the
flux to be orthogonal to $\Omega$ in the Hodge-theoretic sense.

For $K3 \times K3$, the relevant period domain is a product of type IV
domains. Kanno and Watari \cite{KannoWatari2024} parametrize lattice-polarized
K3 surfaces and impose $W=0$ via the Hodge--Riemann bilinear relations on
the flux lattice. Their key result is that CM K3 surfaces --- those whose
transcendental lattice admits an embedding of $\OK$ as endomorphisms --- are
distinguished loci where the moduli are arithmetically constrained.

\subsection{CM K3 surfaces}

\begin{definition}
  A K3 surface $S$ has \emph{complex multiplication by $K = \Q(\sqrt{-D})$}
  if $\End(T(S)) \otimes \Q \supseteq K$, where $T(S)$ is the transcendental
  lattice viewed as a Hodge structure.
\end{definition}

For a CM K3 surface, the period point $[\omega_S] \in \mathbb{P}(T(S) \otimes \C)$
is a CM point in the period domain, meaning it is an algebraic number over
the reflex field of the CM type. This mirrors the theory of CM elliptic
curves, where the $j$-invariant of a CM curve by $K$ is an algebraic integer
in $\HCF(K)$ \cite{Silverman} [CITE-NEEDED: standard reference].

The newform $f$ of level $N$ and weight $k$ associated to a CM K3 surface
via the Shioda--Inose--Morrison correspondence [CITE-NEEDED] encodes the
arithmetic of $T(S)$ in its L-function $L(f,s)$.

\subsection{The $W=0$ locus at CM points}

Kanno and Watari show \cite{KannoWatari2024} that on $K3_1 \times K3_2$
with both factors CM by the same field $K$, the $W=0$ condition reduces to
a system of equations on the flux quanta that is \emph{exactly solvable} at
the CM point. The moduli of both K3 factors are fixed to their CM values,
and the question becomes: how many solutions exist?

The number of solutions to the flux equations, modulo gauge equivalences, is
controlled by the arithmetic of $K$. In the simplest cases, this count is
related to the ideal class group $\mathrm{Cl}(\OK)$ acting on the set of
CM-type flux configurations. We propose that this is the mechanism underlying
Conjecture P1.

%===========================================================================
\section{The Main Conjecture (P1)}
\label{sec:conjecture}
%===========================================================================

Let $K = \Q(\sqrt{-D})$ be an imaginary quadratic field with class number
$h = h(-D)$, ring of integers $\OK$, and Hilbert class field $\HCF(K)$.
Let $S_1, S_2$ be K3 surfaces both of CM type by $K$, with associated
newforms $f_1, f_2$ of weight $5$ [CITE-NEEDED: weight confirmation for
$K3 \times K3$ Hodge structure] and special L-values
$\alpha = L(f_1, 2) / \pi^2$ and $\beta = L(f_2, 2) / \pi^2$ (normalized).

\begin{conjecture}[P1: Class Number One Selection Rule]
\label{conj:P1}
Let $\mathcal{M}_{W=0}(K)$ denote the set of Kanno--Watari $W=0$
flux vacua on $K3_1(K) \times K3_2(K)$ with both K3 factors carrying CM
by $K$. Then:
\begin{enumerate}
  \item[\rm(i)] $|\mathcal{M}_{W=0}(K)| = 1$ \emph{(isolated vacuum)} if and
    only if $h(-D) = 1$.
  \item[\rm(ii)] If $h(-D) = h > 1$, then $|\mathcal{M}_{W=0}(K)| = h$,
    corresponding to the $h$ distinct ideal classes in $\mathrm{Cl}(\OK)$
    acting on CM-type configurations.
  \item[\rm(iii)] In case (i), the vacuum moduli are rational over $K$
    (or $\Q$ after descent). In case (ii), the $h$ vacua are conjugate
    under $\Gal(\HCF(K)/K) \cong \mathrm{Cl}(\OK)$ and none is individually
    rational over $K$.
\end{enumerate}
\end{conjecture}

\begin{remark}
Part (i) is the physically striking claim: the nine Heegner discriminants
$D \in \{1, 2, 3, 7, 11, 19, 43, 67, 163\}$ are precisely those for which
the KW $W=0$ mechanism produces a \emph{single}, arithmetically distinguished
vacuum. For all other CM fields, the string landscape sees a nontrivial
Galois family, not a unique ground state.
\end{remark}

\begin{remark}
The conjecture as stated is \emph{conjectural at the level of string physics}.
Part (iii) has independent mathematical content that is partially proved
(see Section~\ref{sec:evidence}). Parts (i) and (ii) require the additional
physical identification of Galois orbits with distinct flux vacua, which
is the heuristic gap discussed in Section~\ref{sec:gap}.
\end{remark}

%===========================================================================
\section{Evidence: M142 L-Value Rationality Hierarchy}
\label{sec:evidence}
%===========================================================================

The strongest available evidence for Conjecture P1 comes from a proved
hierarchy of rationality properties for special L-values of CM newforms,
which we refer to as the M142 hierarchy (following internal ECI-framework
notation).

\subsection{The M142 Theorem (proved)}

Let $f$ be the weight-$5$ newform associated to a CM K3 surface by
$K = \Q(\sqrt{-D})$, and define
\[
  \alpha_2 = \alpha_2(f, K) \coloneqq \frac{L(f, 2)}{\pi^2 \cdot \Omega_f},
\]
where $\Omega_f$ is a canonical Neron-type period of $f$
[CITE-NEEDED: precise period normalization].

\begin{theorem}[M142 Hierarchy, conditional on Lemma C of the companion paper]
\label{thm:M142}
For each of the nine Heegner discriminants $D \in \{1, 2, 3, 7, 11, 19,
43, 67, 163\}$, and the associated weight-$5$ CM newform $f$,
\[
  \alpha_2(f, K) \in \Q
\]
holds conditionally on the classical Damerell--Shimura rationality at
the right-edge critical integer; an explicit verification at 96--213
digit PARI precision is recorded for each of the nine entries in the
companion paper.
\end{theorem}

This statement combines the classical Damerell--Shimura rationality
theorem at the right-edge critical integer with the specific algebraic
structure of weight-5 newforms at Heegner discriminants. The classical
algebraicity over the Hilbert class field is due to Shimura
\cite{Shimura1977}; the rationality (rather than algebraicity) at
$h_K=1$ follows because $\HK = K$ and the form has rational Fourier
coefficients.

\subsection{The $h > 1$ counterpart}

Classical results in algebraic number theory, dating to Gross--Zagier type
formulas and Goldfeld--Szpiro theory, establish the complementary statement:

\begin{proposition}[Classical; see GS81 \protect{[CITE-NEEDED]}]
\label{prop:h>1}
For $K = \Q(\sqrt{-D})$ with $h(-D) = h > 1$, the normalized L-value
$\alpha_2(f, K)$ is algebraic of degree $h$ over $\Q$, with minimal
polynomial of degree $h$. Its Galois conjugates under $\Gal(\HCF(K)/K)$
are the $h$ distinct normalized L-values corresponding to the $h$ CM types
(ideal class representatives).
\end{proposition}

\begin{example}[$D = 88$, $h = 2$]
For $K = \Q(\sqrt{-88})$, which has class number $2$, the ideal class group
$\mathrm{Cl}(\OK) \cong \Z/2\Z$ acts on the two CM types, and
$\alpha_2 \in \HCF(K) \setminus K$. The two conjugate values
$\alpha_2$ and $\sigma(\alpha_2)$ (where $\sigma$ generates
$\Gal(\HCF(K)/K)$) are not individually rational. This is consistent with
Conjecture P1(ii) predicting two flux vacua for $D=88$.
\end{example}

\begin{remark}[Numerical verification caveat]
Direct numerical verification of $\alpha_2$ via partial Dirichlet sums
$\sum_{n \le N} a_n / n^2$ is invalid: for a weight-5 newform, the
L-function $L(f, s)$ has analytic conductor $s_0 = 3$ (roughly), so $s=2$
lies inside the critical strip where the Dirichlet series does not converge
absolutely. Partial sums oscillate without converging. Rigorous evaluation
requires analytic continuation via the functional equation or PARI/GP's
\texttt{lfun} module. This lesson was learned explicitly during ECI-framework
development (M142 numerical false alarm, 2026-05-08).
\end{remark}

\subsection{Physical interpretation of rationality}

The bridge from L-value rationality to vacuum counting requires the following
chain:
\begin{enumerate}
  \item The KW flux equations at a CM point are governed by period integrals
    of the CM K3.
  \item Period integrals at CM points are (up to transcendental factors)
    equal to $\alpha_2$ and its algebraic conjugates.
  \item A ``vacuum'' corresponds to a simultaneous solution; when $\alpha_2 \in \Q$,
    all algebraic conjugates coincide, giving one solution.
  \item When $\alpha_2 \notin \Q$, the $h(-D)$ conjugates are distinct, giving
    $h(-D)$ physically inequivalent solutions.
\end{enumerate}
Step 2 is the most heuristic link, discussed further in Section~\ref{sec:gap}.

%===========================================================================
\section{Hilbert Class Field Interpretation}
\label{sec:hcf}
%===========================================================================

The Hilbert class field $\HCF(K)$ is the maximal unramified abelian extension
of $K$. By class field theory,
\[
  \Gal(\HCF(K)/K) \cong \mathrm{Cl}(\OK).
\]
For $h(-D) = 1$, this is trivial: $\HCF(K) = K$ and every ideal in $\OK$
is principal.

\subsection{CM values and HCF}

The theory of complex multiplication for elliptic curves gives:

\begin{theorem}[Kronecker, Weber, Takagi--Shimura; see Silverman \protect{[CITE-NEEDED]}]
If $E$ is an elliptic curve with CM by $\OK$, then the $j$-invariant
$j(E) \in \HCF(K)$ and generates $\HCF(K)$ over $K$.
\end{theorem}

The analogous statement for K3 surfaces with CM is expected but not fully
proved in general:

\begin{conjecture}[CM values for K3, expected; \protect{[CITE-NEEDED]}]
If $S$ is a K3 surface with CM by $\OK$ of rank 2 (algebraic K3), then
the relevant period ratios lie in $\HCF(K)$ and their $\Gal(\HCF(K)/K)$-orbit
has cardinality $h(-D)$.
\end{conjecture}

For our purposes, we use the $h=1$ case where these expected results are
unconditional: when $K$ is one of the nine Heegner fields, $\HCF(K) = K$,
so all CM values are already rational over $K$ (and often over $\Q$ after
the trace map).

\subsection{Galois action on vacua}

If we accept (provisionally) that the KW flux vacua at a CM point form a
set $\mathcal{M}_{W=0}(K)$ on which $\mathrm{Cl}(\OK)$ acts freely and
transitively, then:
\[
  |\mathcal{M}_{W=0}(K)| = |\mathrm{Cl}(\OK)| = h(-D).
\]
The $h(-D) = 1$ case gives the unique vacuum predicted by Conjecture P1(i).

This is precisely the string-theory analogue of the statement that a CM
elliptic curve over $\Q$ with CM by $\OK$ of class number 1 has a rational
model, while for $h > 1$ one needs to pass to the class field to find the
model.

%===========================================================================
\section{Heuristic vs.\ Rigorous Gap}
\label{sec:gap}
%===========================================================================

We are explicit about what is proved versus conjectural.

\subsection{What is proved}

\begin{enumerate}
  \item Theorem~\ref{thm:M142} (M142): $\alpha_2(f,K) \in \Q$ for all nine
    Heegner fields. This is a rigorous mathematical statement.
  \item Proposition~\ref{prop:h>1}: For $h > 1$, $\alpha_2(f,K) \notin \Q$
    and lies in $\HCF(K)$. This follows from classical CM theory applied
    to the appropriate automorphic L-function (pending rigorous reference
    for weight-5 case; [CITE-NEEDED]).
  \item The nine Heegner discriminants are the complete list of imaginary
    quadratic fields with $h=1$ (Baker--Stark theorem \cite{Baker1966,Stark1967}
    [CITE-NEEDED: precise Baker and Stark references]).
\end{enumerate}

\subsection{What is heuristic}

\begin{enumerate}
  \item \textbf{The period--L-value bridge}: That the KW flux equations at
    a CM point are controlled by $\alpha_2(f,K)$ specifically requires
    identifying the relevant periods with L-values via the Shimura period
    relations. This identification is natural but has not been written out
    rigorously in the KW context.

  \item \textbf{Galois action on flux vacua}: That the $h(-D)$ algebraic
    conjugates of $\alpha_2$ correspond bijectively to $h(-D)$ physically
    distinct flux vacua requires a more careful analysis of the KW moduli
    equations. Different ideal class representatives could in principle
    give gauge-equivalent vacua.

  \item \textbf{UV stability}: The conjecture includes UV stability of the
    vacuum under string corrections. While supersymmetric $W=0$ vacua are
    typically protected, a complete argument would require specifying the
    K\"ahler potential and verifying no-scale structure.

  \item \textbf{Weight-5 vs.\ weight-2}: The M142 result is stated for
    weight-5 newforms. The precise weight of the newform attached to the
    transcendental lattice of a K3 surface depends on the Hodge numbers;
    for a generic K3 the relevant weight in the Langlands correspondence
    may differ. This requires case-by-case verification.
\end{enumerate}

\subsection{Falsifiability}

Conjecture P1 is falsifiable:
\begin{itemize}
  \item If a CM field $K$ with $h(-D) = 1$ admits more than one KW flux
    vacuum, Conjecture P1(i) fails.
  \item If a CM field $K$ with $h(-D) = 2$ admits exactly one KW flux
    vacuum, Conjecture P1(ii) fails.
  \item If $\alpha_2(f,K)$ is found to be irrational for some Heegner
    discriminant (via rigorous PARI \texttt{lfun} computation), Theorem~\ref{thm:M142}
    is falsified and the entire chain collapses.
\end{itemize}

%===========================================================================
\section{Open Problems and Computational Tests}
\label{sec:open}
%===========================================================================

\subsection{Computational tests (near-term)}

\begin{enumerate}
  \item \textbf{PARI/GP verification of $\alpha_2$ for all nine Heegner fields.}
    Use \texttt{lfun(f, 2)} for each associated weight-5 CM newform from
    LMFDB, and verify $\alpha_2 \in \Q$ to 50+ decimal digits. This is the
    most immediate check and is feasible on current hardware within hours.

  \item \textbf{$D=88$ case study.}
    For $K = \Q(\sqrt{-88})$ ($h=2$), compute both L-values corresponding
    to the two ideal classes, verify they are quadratic conjugates over $K$,
    and confirm their minimal polynomial has degree 2.

  \item \textbf{Explicit KW equation count.}
    For a concrete lattice-polarized K3 of CM type by $\Q(\sqrt{-3})$
    ($h=1$) versus $\Q(\sqrt{-5})$ ($h=2$), enumerate solutions to the
    KW flux orthogonality conditions and verify the predicted counts 1 vs.\ 2.

  \item \textbf{LMFDB newform identification.}
    For each Heegner discriminant, identify the associated CM newforms in
    LMFDB \cite{LMFDB} [CITE-NEEDED: LMFDB standard citation] and record
    their levels and character data.
\end{enumerate}

\subsection{Theoretical open problems}

\begin{enumerate}
  \item \textbf{Rigorous period--L-value bridge.}
    Extend Shimura's period relations to the KW setting, making precise the
    identification of flux vacuum moduli with $\alpha_2$.

  \item \textbf{Galois action on flux vacua.}
    Prove or disprove that $\mathrm{Cl}(\OK)$ acts freely and transitively
    on $\mathcal{M}_{W=0}(K)$.

  \item \textbf{Non-CM generalization.}
    Does a weaker form of P1 hold for K3 surfaces with real multiplication
    or other arithmetic structures? The class number 1 condition would need
    reinterpretation.

  \item \textbf{Connection to Swampland.}
    Does Conjecture P1 have implications for swampland conjectures? An
    isolated vacuum ($h=1$) is UV-stable and sits at a definite point in
    the landscape; a Galois family ($h>1$) forms a discrete but non-isolated
    set of points. The de Sitter conjecture and distance conjecture
    [CITE-NEEDED] could constrain which regime is physical.

  \item \textbf{$F$-theory lift.}
    The $K3 \times K3$ geometry admits an $F$-theory realization. Does the
    P1 selection rule have a natural $F$-theory interpretation in terms of
    elliptic fibration arithmetic?

  \item \textbf{Arithmetic of $\alpha_2$ and height bounds.}
    For $h > 1$, can one bound the height of $\alpha_2$ as an element of
    $\HCF(K)$ in terms of $D$? Such a bound would give quantitative
    control over how ``close'' the $h(-D)$ vacua are in moduli space.
\end{enumerate}

%===========================================================================
\section*{Acknowledgements}
%===========================================================================

[To be completed.] We thank [collaborators] for discussions. The ECI
framework computations used PARI/GP and LMFDB data.

%===========================================================================
% Bibliography
%===========================================================================

\begin{thebibliography}{99}

\bibitem{KannoWatari2024}
S.~Kanno and T.~Watari,
\textit{$W=0$ Flux Vacua on K3$\times$K3 and Arithmetic},
arXiv:2012.01111 [hep-th] (2020). [Verified]

\bibitem{Moore1998}
G.~Moore,
\textit{Arithmetic and Attractors},
arXiv:hep-th/9807087 (1998). [Verified]

\bibitem{BIKMSV1996}
M.~Bershadsky, K.~Intriligator, S.~Kachru, D.~Morrison, V.~Sadov, and C.~Vafa,
\textit{Geometric Singularities and Enhanced Gauge Symmetries},
Nucl.\ Phys.\ \textbf{B481} (1996) 215--252,
arXiv:hep-th/9605200. [Verified]

\bibitem{Shimura1977}
G.~Shimura,
\textit{On the periods of modular forms},
Math.\ Ann.\ \textbf{229} (1977) 211--221. [CITE-NEEDED: verify page range and volume]

\bibitem{Baker1966}
A.~Baker,
\textit{Linear forms in the logarithms of algebraic numbers},
Mathematika \textbf{13} (1966) 204--216. [CITE-NEEDED: verify]

\bibitem{Stark1967}
H.~M.~Stark,
\textit{A complete determination of the complex quadratic fields of class-number one},
Michigan Math.\ J.\ \textbf{14} (1967) 1--27. [CITE-NEEDED: verify]

\bibitem{Silverman}
J.~H.~Silverman,
\textit{Advanced Topics in the Arithmetic of Elliptic Curves},
Springer GTM 151, 1994. [CITE-NEEDED: standard reference]

\bibitem{LMFDB}
The LMFDB Collaboration,
\textit{The L-functions and Modular Forms Database},
\url{https://www.lmfdb.org} (2023). [CITE-NEEDED: use current citation]

\end{thebibliography}

\end{document}
